% !TeX program = tectonic
\documentclass[UTF8,a4paper,11pt]{ctexart}

\usepackage[left=2.45cm,right=2.45cm,top=2.35cm,bottom=2.35cm,headheight=15pt]{geometry}
\usepackage{amsmath,amssymb}
\usepackage{array,booktabs,tabularx}
\usepackage{caption}
\usepackage{fancyhdr}
\usepackage{float}
\usepackage{listings}
\usepackage{tikz}
\usepackage{xcolor}
\usepackage{xurl}
\usepackage{hyperref}
\usetikzlibrary{arrows.meta,positioning}

\hypersetup{
  colorlinks=true,
  linkcolor=blue!45!black,
  citecolor=blue!45!black,
  urlcolor=blue!55!black,
  pdftitle={基于 Clang/LLVM 的语言处理流程与 SysY 手写实现},
  pdfauthor={LIZZYGREAT},
  pdfsubject={编译原理实验报告}
}
\urlstyle{same}
\setlength{\parindent}{2em}
\setlength{\parskip}{0.25em}
\setlength{\emergencystretch}{2.5em}
\tolerance=2000
\hbadness=2500
\renewcommand{\arraystretch}{1.22}
\newcolumntype{Y}{>{\raggedright\arraybackslash}X}
\newcolumntype{P}[1]{>{\raggedright\arraybackslash}p{#1}}
\newcommand{\code}[1]{\texttt{\detokenize{#1}}}
\captionsetup{font=small,labelfont=bf,skip=5pt}

\pagestyle{fancy}
\fancyhf{}
\fancyhead[L]{\small 编译原理第一次作业实验报告}
\fancyhead[R]{\small Clang/LLVM 与 SysY}
\fancyfoot[C]{\thepage}
\renewcommand{\headrulewidth}{0.4pt}
\renewcommand{\footrulewidth}{0pt}

\lstset{
  basicstyle=\ttfamily\small,
  backgroundcolor=\color{black!3},
  frame=single,
  rulecolor=\color{black!15},
  framerule=0.4pt,
  xleftmargin=1em,
  xrightmargin=1em,
  columns=fullflexible,
  keepspaces=true,
  showstringspaces=false,
  breaklines=true,
  aboveskip=0.6em,
  belowskip=0.7em
}

\title{\heiti 基于 Clang/LLVM 的语言处理流程与 SysY 手写实现}
\author{作者：LIZZYGREAT}
\date{2026 年 10 月 7 日}

\begin{document}
\maketitle

\begin{abstract}
本文以一个含函数、while、if/else、取模和整数输入输出的 SysY 风格程序为贯穿案例，使用 Clang 14 观察预处理、词法/语法/语义分析与 IR 生成、x86-64 汇编、ELF 目标文件和显式链接的完整过程，并通过实际 \code{readelf} 输出核对节区、符号与重定位项。随后根据源程序语义独立手写 LLVM IR 与 RV64 汇编，对照控制流图与栈帧布局逐条说明设计意图，并分别链接运行库验证。六组输入在宿主机 reference、手写 LLVM 和 QEMU/RISC-V 上均得到 1、5、8、16、21、85。优化对比显示，Clang -O2 消除了该模块中的 alloca/load/store，并在 x86-64 上生成四 lane 向量运算；RV64 目标保留标量循环并内联 calc。进阶部分从源码构建了 AscendNPU-IR 的 bishengir-opt，逐字执行上游自带的 RUN 命令并用其 CHECK 断言做 FileCheck，PlanMemory 与 InjectSync 两组用例均通过，并保存了真实 Before/After；未安装 CANN 或使用 NPU 设备，该范围不属本任务所需。全部实验由脚本驱动，可逐阶段复现。

\noindent\textbf{关键词：} 编译器；Clang；LLVM IR；RISC-V；ABI；MLIR；lowering
\end{abstract}

\section{引言}

编译器把源程序逐步转换为可执行程序。理解每一步的输入、输出和符号关系，有助于分清预处理、编译、汇编和链接的职责。LLVM IR 适合观察语言结构如何变为控制流和算术指令；RISC-V 汇编则展示调用约定、寄存器和栈帧如何落实到目标架构。

本实验固定一个程序，避免不同算法导致比较基准变化。程序令 $i$ 从 1 递增到输入 $n$，若 $i$ 为偶数则将 $2i$ 加入 sum，否则将 $i$ 加入 sum，最后输出 sum。程序覆盖局部整型变量、赋值、加法、乘法、取模、比较、条件分支、循环、函数参数、函数返回和运行库调用。所有实现都使用同一输入集合验证。

全文依次回答四个问题：预处理、编译、汇编、链接各自做了什么，产物之间是什么关系（第 3 节）；源程序的条件与循环如何变成 LLVM 基本块和 \code{br}，又如何独立手写出一份等价 IR（第 4 节）；同一逻辑在 RV64 上如何落到寄存器、栈帧和调用约定（第 5 节）；优化级别和抽象层次的变化会带来什么（第 7、8 节）。所有结论都来自本机真实产物，原始命令、输出和中间文件分别保存在 \code{scripts/}、\code{logs/} 和 \code{records/} 中，第 9 节给出复现入口。

\section{实验对象与工具环境}

\subsection{贯穿程序}

整次作业只使用一个程序，避免更换算法导致比较基准漂移。其 SysY 源程序 \code{source/example.sy} 如下。

\begin{lstlisting}
int calc(int n) {
    int i = 1;
    int sum = 0;
    while (i <= n) {
        if (i % 2 == 0) { sum = sum + i * 2; }
        else            { sum = sum + i;     }
        i = i + 1;
    }
    return sum;
}

int main() {
    int n;      n = getint();
    int result; result = calc(n);
    putint(result);
    return 0;
}
\end{lstlisting}

该程序覆盖局部整型变量、赋值、整型常量、加法、乘法、取模、比较、\code{if/else}、\code{while}、函数定义与调用、参数传递、返回值以及运行库输入输出，足以观察编译器的主要结构。语义为：$i$ 从 1 递增到 $n$，$i$ 为偶数时累加 $2i$，否则累加 $i$，最后输出 sum。为使 Clang 能直接处理，\code{source/example_clang.c} 与 \code{example.sy} 算法完全一致，仅补充运行库声明 \code{int getint(void);} 与 \code{void putint(int x);}；宿主机验证所需的 \code{getint}/\code{putint} 由 \code{source/host_runtime.c} 以 \code{scanf}/\code{printf} 实现。固定测试集为 $n=1,2,3,4,5,10$，期望输出依次为 $1,5,8,16,21,85$；该期望值独立于任何实现，作为唯一的 Ground Truth。

\subsection{工具环境}

工作环境为 Ubuntu 22.04.5 on WSL2（内核 \code{6.18.33.2-microsoft-standard-WSL2}），全部工具版本记录于 \code{logs/00_environment.log}。

\begin{table}[H]
  \centering\small
  \caption{本实验使用的工具链}
  \begin{tabularx}{\textwidth}{P{4.1cm}P{4.0cm}Y}
    \toprule
    工具 & 版本 & 本实验用途 \\
    \midrule
    Clang / LLVM & 14.0.0 & 预处理、LLVM IR、x86-64 汇编、目标文件；接受手写 \code{.ll} \\
    GCC & 11.4.0 & 提供系统链接器驱动环境 \\
    RISC-V unknown-elf GCC & 15.1.0 & 汇编、链接手写 RV64 汇编与课程 runtime \\
    QEMU & 7.0.0 & 用户态运行 RISC-V 可执行文件 \\
    Flex / Bison & 2.6.4 / 3.8.2 & 课程编译器环境既有组件 \\
    Tectonic & 0.16.9 & 编译本报告 LaTeX 源文件 \\
    \bottomrule
  \end{tabularx}
\end{table}

课程目录 \code{SysY_Compiler_2025} 已提供编译器与 SysY runtime，本实验只复用其 \code{lib/libsysy_riscv.a}，没有重新构建或改动该目录。Clang 的 RISC-V 参考目标与手写目标统一使用 \code{rv64gc/lp64d}，与 runtime 的 ELF 架构标志一致。

\section{Clang 语言处理流程}

\begin{figure}[H]
  \centering
  \begin{tikzpicture}[
    stage/.style={draw=blue!45!black,fill=blue!4,rounded corners=2pt,
      align=center,text width=1.82cm,minimum height=1.05cm,font=\small},
    flow/.style={-{Stealth[length=2mm]},thick,draw=blue!55!black},
    node distance=0.48cm
  ]
    \node[stage] (src) {SysY/C\\源文件};
    \node[stage,right=of src] (pre) {预处理\\\code{.i}};
    \node[stage,right=of pre] (ir) {LLVM IR\\\code{.ll}};
    \node[stage,right=of ir] (asm) {汇编/目标\\\code{.s/.o}};
    \node[stage,right=of asm] (elf) {显式链接\\ELF};
    \node[stage,right=of elf] (run) {执行\\host\\QEMU};
    \draw[flow] (src) -- (pre);
    \draw[flow] (pre) -- (ir);
    \draw[flow] (ir) -- (asm);
    \draw[flow] (asm) -- (elf);
    \draw[flow] (elf) -- (run);
  \end{tikzpicture}
  \caption{从源程序到运行文件的实验流程}
  \label{fig:pipeline}
\end{figure}

\subsection{预处理}

命令 \code{clang -E source/example_clang.c -o generated/example.i} 生成预处理文件。预处理器只处理以 \code{#} 开头的指令：展开宏、插入 \code{#include} 内容、求值条件编译、处理行标记。它按行替换文本，不做词法之外的语义判断，因此不理会 \code{while}、\code{if} 的含义。本例没有宏、条件编译和头文件，\code{generated/example.i} 只有 11 行 \code{#} 行标记，其后语句与输入完全相同。

\begin{lstlisting}
# 1 "source/example_clang.c"
# 1 "<built-in>" 1
# 1 "<built-in>" 3
# 361 "<built-in>" 3
...
# 1 "source/example_clang.c" 2
int getint(void);
void putint(int x);

int calc(int n) {
    int i = 1;
    ...
\end{lstlisting}

行标记形如 \code{# 行号 "文件名" 标志}，供后续阶段和调试器把生成代码映射回源文件；\code{<built-in>} 是编译器内建的定义集合。除这 11 行标记外，\code{.i} 内容与输入逐字相同。预处理器不建立符号表、不检查类型、也不理解控制流，这些工作都在下一节的编译阶段完成。

\subsection{编译器内部各阶段}

\code{clang -cc1} 是真正的编译器前端。对同一份源文件，它依次完成以下工作；本实验通过在 \code{example.ll} 中看到的结果反推各阶段产物。

\begin{table}[H]
  \centering\small
  \caption{编译器内部阶段及其在本程序中的表现}
  \begin{tabularx}{\textwidth}{P{2.5cm}P{4.6cm}Y}
    \toprule
    阶段 & 职责 & 本程序的观察 \\
    \midrule
    词法分析 & 把字符流切成记号 & \code{int}、\code{calc}、\code{i}、\code{<=}、\code{100} 等记号 \\
    语法分析 & 按文法构造抽象语法树 & 得到 \code{while}、\code{if/else} 的嵌套结构 \\
    语义分析 & 类型检查、建立符号表、插入隐式转换 & 判定 \code{i <= n} 为 \code{i32} 比较，\code{calc(n)} 类型匹配 \\
    IR 生成 & 遍历 AST 生成 LLVM IR & 局部变量变为 \code{alloca}，语句变为基本块与 \code{br} \\
    中端优化 & 在 IR 上做与目标无关的变换 & \code{-O0} 不做；\code{-O2} 提升为 SSA 与 \code{phi}（第 7 节） \\
    后端 & 指令选择、寄存器分配、指令调度 & 生成 x86-64 或 RV64 汇编 \\
    \bottomrule
  \end{tabularx}
\end{table}

其中前三个阶段常合称“前端”，它们只依赖源语言；中端依赖 IR；后端依赖目标架构。这与图~\ref{fig:pipeline} 的整体流程互相印证：预处理器独立于前端，前端到 IR 为止，中端和后端才与目标相关。

\subsection{LLVM IR 生成}

命令如下：
\begin{lstlisting}
clang -S -emit-llvm -O0 source/example_clang.c -o generated/example.ll
\end{lstlisting}
生成 \code{generated/example.ll}。它是与目标无关的中间表示，函数、类型、指令都在此层显式可见。\code{-O0} 下一个 \code{alloca} 对应一个局部变量，每次读写都显式 \code{load}/\code{store}，见第 4 节。

\subsection{汇编与目标文件}

\code{clang -S -O0} 生成宿主机汇编 \code{generated/example_x86.s}；\code{clang -c -O0} 继续调用集成汇编器，将汇编编码为机器码，产出 ELF64 可重定位目标文件 \code{generated/example.o}。汇编器负责指令编码、地址分配和符号表构造，此时地址尚未最终确定，因此外部引用以重定位项的形式保留。对 \code{example.o} 的实际检查结果如下。

\begin{table}[H]
  \centering\small
  \caption{\code{example.o} 的 ELF 特征（\code{file} 与 \code{readelf} 实际输出）}
  \begin{tabularx}{\textwidth}{P{3.0cm}Y}
    \toprule
    检查项 & 实际结果 \\
    \midrule
    文件类型 & \code{ELF 64-bit LSB relocatable, x86-64, not stripped}，即 \code{ET_REL} \\
    主要节区 & \code{.text}（162 字节）、\code{.rela.text}、\code{.eh_frame}、\code{.symtab}、\code{.strtab} \\
    已定义符号 & \code{calc}（\code{FUNC}，100 字节）、\code{main}（\code{FUNC}，50 字节） \\
    未定义符号 & \code{getint}、\code{putint}，均为 \code{UND} \\
    重定位项 & \code{.rela.text} 中 3 条 \code{R_X86_64_PLT32}，目标为 \code{getint}、\code{calc}、\code{putint} \\
    \bottomrule
  \end{tabularx}
\end{table}

\code{readelf -h} 给出的 ELF 文件头也印证了这一点：

\begin{lstlisting}
Class:                     ELF64
Type:                      REL (Relocatable file)
Machine:                   Advanced Micro Devices X86-64
Entry point address:       0x0
Start of program headers:  0 (bytes into file)
\end{lstlisting}

\code{Type} 为 \code{REL} 而非 \code{EXEC}/\code{DYN}，\code{Entry point} 为 \code{0x0}，且程序头数量为 0，说明该文件只是“待合并的部件”：地址未定、没有进程映像布局，无法被内核直接加载执行。符号表与重定位项还提供了两点信息。其一，符号表把 \code{calc}/\code{main} 标为 \code{FUNC} 且带字节大小，而 \code{getint}/\code{putint} 标为 \code{UND}，即本文件提供什么、还欠什么。其二，重定位项中除外部符号外还出现了 \code{calc}——本机 Clang 默认使用 PIE，模块内调用同样经 PLT 重定位，其最终地址由链接器确定，这也解释了为什么 \code{.rela.text} 中有 3 条记录而不是 2 条。

\subsection{链接}

本实验分别生成 \code{example.o} 与 \code{host_runtime.o}，再用 Clang 驱动显式链接。日志显示 Clang 调用 \code{/usr/bin/ld}，并加入 CRT 启动对象和 \code{-lc}。链接器按符号名把各目标文件的定义与引用配对：\code{host_runtime.o} 定义 \code{getint}/\code{putint}，解决 \code{example.o} 的两个未定义符号；系统 C 库解析 runtime 中的 \code{scanf}/\code{printf}。随后合并同名节区、按最终地址应用重定位、丢弃不再需要的符号信息，生成可执行 ELF \code{generated/reference_linked}，其六组输出与 Ground Truth 一致。反过来说，单独的 \code{example.o} 既缺少两个外部函数定义，也没有进程启动代码，因此不能独立完成整个程序功能；\code{getint}/\code{putint} 正是在链接阶段被解析的。

\begin{table}[H]
  \centering\small
  \caption{语言处理阶段及本实验观察}
  \begin{tabularx}{\textwidth}{P{2.4cm}P{5.5cm}Y}
    \toprule
    阶段 & 命令或产物 & 本实验观察 \\
    \midrule
    预处理 & \code{clang -E}；\code{example.i} & 保留 C 语句并带行标记；本例无宏/头文件展开 \\
    LLVM IR & \code{clang -S -emit-llvm -O0}；\code{example.ll} & alloca、load/store、icmp、br、call、ret \\
    汇编 & \code{clang -S -O0}；\code{example_x86.s} & x86-64 指令和调用序列 \\
    汇编器 & \code{clang -c}；\code{example.o} & ELF REL、节区、符号和重定位 \\
    链接器 & \code{example.o + host_runtime.o} & 解析 runtime/C 库符号并生成可执行 ELF \\
    \bottomrule
  \end{tabularx}
\end{table}

\section{LLVM IR 控制流与手写实现}

\begin{figure}[H]
  \centering
  \begin{tikzpicture}[
    block/.style={draw=blue!45!black,fill=blue!4,rounded corners=2pt,
      align=center,text width=3.15cm,minimum height=0.92cm,font=\small},
    edge/.style={-{Stealth[length=2mm]},thick,draw=blue!55!black}
  ]
    \node[block] (entry) at (0,0) {entry\\初始化 $i$、sum};
    \node[block] (cond) at (0,-1.55) {while 条件\\$i \le n$};
    \node[block] (exit) at (5,-1.55) {while.end\\return sum};
    \node[block] (ifcond) at (0,-3.25) {if 条件\\$i \bmod 2 = 0$};
    \node[block] (then) at (-2,-5.0) {if.then\\sum $\mathrel{+}=2i$};
    \node[block] (else) at (2,-5.0) {if.else\\sum $\mathrel{+}=i$};
    \node[block] (merge) at (0,-6.8) {if.end\\$i{+}{+}$};
    \draw[edge] (entry) -- (cond);
    \draw[edge] (cond.east) -- node[above,font=\footnotesize]{$i>n$} (exit.west);
    \draw[edge] (cond) -- node[right,font=\footnotesize]{继续} (ifcond);
    \draw[edge] (ifcond.west) -- node[above left,font=\footnotesize]{偶数} (then.north);
    \draw[edge] (ifcond.east) -- node[above right,font=\footnotesize]{奇数} (else.north);
    \draw[edge] (then.south) -- (merge.north west);
    \draw[edge] (else.south) -- (merge.north east);
    \draw[edge] (merge.west) -- (-2.55,-6.8) -- (-2.55,-1.55) -- (cond.west);
  \end{tikzpicture}
  \caption{\code{calc} 的基本块与循环回边}
\end{figure}

Clang -O0 生成的 calc 使用数字标签。入口块初始化 i 和 sum 后进入 \texttt{\%5}；\texttt{\%5} 用 \code{icmp sle} 计算 $i \le n$。真分支进入 \texttt{\%9} 判断 $i\bmod 2$，两个分支 \texttt{\%13}、\texttt{\%18} 分别计算偶数项和奇数项，并汇入 \texttt{\%22}。\texttt{\%22} 执行 $i{+}{+}$ 后无条件跳回 \texttt{\%5}，构成 while 回边；假分支到 \texttt{\%25} 读取 sum 并返回。

\begin{table}[H]
  \centering\small
  \caption{源程序结构与 LLVM 指令对应}
  \begin{tabularx}{\textwidth}{P{3.2cm}P{3.1cm}Y}
    \toprule
    源程序结构 & LLVM 指令 & 作用 \\
    \midrule
    局部变量声明/赋值 & alloca、store & 为 n/i/sum 分配栈槽并写初值或新值 \\
    读取变量 & load & 取 i、n、sum 的当前值 \\
    加法、乘法、取模 & add、mul、srem & 计算计数器、累加项与奇偶性 \\
    比较 & icmp sle、icmp eq & 判断循环范围和余数是否为零 \\
    if/while & br、基本块 & 根据条件转移，合并分支并形成循环回边 \\
    函数关系 & define、call、ret & 定义 calc/main，调用 getint/calc/putint 并返回 \\
    \bottomrule
  \end{tabularx}
\end{table}

\subsection{手写 LLVM IR}

手写 \code{llvm/manual.ll} 不复用 Clang 的输出，而是按上面的 CFG 直接组织命名基本块。与自动生成版本相比，手写版做了几处不同选择：块名用语义化的 \code{while.cond} 而不是 \texttt{\%5}；不产生 \code{attributes}、\code{!llvm.loop} 等元数据；\code{main} 中直接用 SSA 值传递 \code{getint} 与 \code{calc} 的返回值，不再绕经栈槽。控制流与算术指令仍与源程序逐项对应。

\begin{lstlisting}
declare i32 @getint()
declare void @putint(i32)

define i32 @calc(i32 %n) {
entry:
  %i.addr = alloca i32, align 4
  %sum.addr = alloca i32, align 4
  store i32 1, i32* %i.addr, align 4      ; i = 1
  store i32 0, i32* %sum.addr, align 4    ; sum = 0
  br label %while.cond

while.cond:
  %i.cond = load i32, i32* %i.addr, align 4
  %continue = icmp sle i32 %i.cond, %n    ; i <= n
  br i1 %continue, label %if.cond, label %while.end

if.cond:
  %i.parity = load i32, i32* %i.addr, align 4
  %remainder = srem i32 %i.parity, 2      ; i % 2
  %is.even = icmp eq i32 %remainder, 0
  br i1 %is.even, label %if.then, label %if.else

if.then:
  %sum.even = load i32, i32* %sum.addr, align 4
  %i.even = load i32, i32* %i.addr, align 4
  %twice.i = mul nsw i32 %i.even, 2       ; 2 * i
  %sum.even.next = add nsw i32 %sum.even, %twice.i
  store i32 %sum.even.next, i32* %sum.addr, align 4
  br label %if.end

if.else:
  %sum.odd = load i32, i32* %sum.addr, align 4
  %i.odd = load i32, i32* %i.addr, align 4
  %sum.odd.next = add nsw i32 %sum.odd, %i.odd
  store i32 %sum.odd.next, i32* %sum.addr, align 4
  br label %if.end

if.end:
  %i.current = load i32, i32* %i.addr, align 4
  %i.next = add nsw i32 %i.current, 1     ; i = i + 1
  store i32 %i.next, i32* %i.addr, align 4
  br label %while.cond                    ; while 的回边

while.end:
  %sum.result = load i32, i32* %sum.addr, align 4
  ret i32 %sum.result
}

define i32 @main() {
entry:
  %n = call i32 @getint()
  %result = call i32 @calc(i32 %n)
  call void @putint(i32 %result)
  ret i32 0
}
\end{lstlisting}

\code{while} 没有也不需要有专门的 LLVM 指令：\code{while.cond} 计算条件、\code{br i1} 二选一，\code{if.end} 的无条件 \code{br} 跳回 \code{while.cond} 即构成回边。\code{if/else} 同样只是 \code{if.cond} 处的一次条件 \code{br} 加两个分支块，两条分支在 \code{if.end} 汇合。Clang 接受该 IR，将其与 \code{host_runtime.c} 链接为 \code{generated/manual_llvm}，六组输入全部通过（\code{outputs/manual-llvm-test.txt}）。

\begin{table}[H]
  \centering\small
  \caption{自动生成 IR 与手写 IR 的组织方式对比}
  \label{tab:ir-compare}
  \begin{tabularx}{\textwidth}{P{2.7cm}P{5.3cm}Y}
    \toprule
    对比项 & \code{generated/example.ll}（Clang -O0） & \code{llvm/manual.ll}（手写） \\
    \midrule
    基本块命名 & 数字标签 \texttt{\%5}、\texttt{\%9}、\texttt{\%13} & 语义名 \code{while.cond}、\code{if.then}、\code{if.end} \\
    局部变量 & 三个 \code{alloca}，含形参 \code{n} 的栈槽 & 只给 \code{i}、\code{sum} 两个 \code{alloca}，\code{n} 直接用 SSA 形参 \\
    分支合流 & \texttt{\%13}、\texttt{\%18} 汇入 \texttt{\%22} & \code{if.then}、\code{if.else} 汇入 \code{if.end} \\
    元数据 & 带 \code{attributes}、\code{!llvm.loop}、\code{!llvm.module.flags} & 不产生任何元数据 \\
    目标三元组 & \code{x86\_64-pc-linux-gnu} 与完整 \code{datalayout} & 仅声明 \code{target triple} \\
    指令集合 & 同一组核心指令 & 逐条对应源语句，无新增或缺失的操作 \\
    \bottomrule
  \end{tabularx}
\end{table}

两份 IR 使用的指令集合完全一致：\code{alloca}、\code{load}、\code{store}、\code{add}、\code{mul}、\code{srem}、\code{icmp}、\code{br}、\code{call}、\code{ret} 逐条对应，没有新增或缺失的操作（表~\ref{tab:ir-compare}）。可见 \code{-O0} 下 LLVM IR 到源程序的映射比较直接，两份实现的差别只在命名、元数据和值表示的组织方式。

\section{RISC-V 汇编与 ABI}

\subsection{参考汇编与调用约定}

RISC-V 参考汇编由 Clang 14 交叉生成（\code{--target=riscv64-unknown-elf -march=rv64gc -mabi=lp64d}），先用于观察成熟后端如何表达同一逻辑，再作为手写实现的对照。其 \code{calc} 序言为：

\begin{lstlisting}
calc:
    addi sp, sp, -32
    sd   ra, 24(sp)
    sd   s0, 16(sp)
    addi s0, sp, 32
    sw   a0, -20(s0)
    ...
.LBB0_1:
    lw   a1, -24(s0)      # i
    lw   a0, -20(s0)      # n
    blt  a0, a1, .LBB0_6
    ...
.LBB0_2:
    lwu  a0, -24(s0)      # i % 2 的替代实现
    srli a1, a0, 31
    add  a1, a1, a0
    andi a1, a1, -2
    subw a0, a0, a1
    bne  a0, zero, .LBB0_4
\end{lstlisting}

与下一节的手写版本对照，有三处差异。其一，\emph{栈帧完全一致}：Clang 同样用 32 字节帧、\code{ra} 存 \code{24(sp)}、\code{s0} 存 \code{16(sp)}、三个局部量放在 \code{-20/-24/-28(s0)}，说明手写实现在没有照抄的情况下收敛到了与真实后端相同的布局。其二，\emph{取模的实现不同}：RV64 没有硬件取余指令，\code{remw} 需要微码或库调用，因此 Clang 把 $i \bmod 2$ 展开成 \code{lwu/srli/add/andi/subw} 五条定点运算；手写版本直接用 \code{remw} 加 \code{beqz}，更贴近源程序语义但指令条数更少，这也说明“指令少”不等于“更快”。其三，\emph{\code{main} 的帧大小不同}：Clang 的 \code{main} 用 48 字节，因为它把返回值先 \code{sd} 溢出到栈再取出；手写 \code{main} 只用 16 字节，返回值始终留在 \code{a0} 中直接传递。

\begin{table}[H]
  \centering\small
  \caption{本实验用到的 RISC-V 寄存器与调用约定}
  \begin{tabularx}{\textwidth}{P{3.0cm}Y}
    \toprule
    寄存器/规则 & 本实验用法 \\
    \midrule
    \code{a0} & 第一个整型参数与整型返回值；getint、calc、putint 之间传值 \\
    \code{t0-t2} & calc 内的 caller-saved 临时值；本实验只用这三个 \\
    \code{s0} & calc 的 frame pointer；属被调用者保存，成对保存/恢复 \\
    \code{sp} & 栈指针；calc/main 分别保留 32/16 字节且保持 16 字节对齐 \\
    \code{ra} & 返回地址；calc 与 main 跨越函数调用保存/恢复 \\
    \bottomrule
  \end{tabularx}
\end{table}

\subsection{栈帧布局}

RV64 调用约定要求 \code{sp} 16 字节对齐。\code{calc} 的栈帧共 32 字节：顶部保存返回地址与旧的 \code{s0}，中部放三个局部变量。堆栈从高地址向低地址增长，因此“入栈”即减小 \code{sp}。

\begin{figure}[H]
  \centering
  \begin{tikzpicture}[font=\small,x=1cm,y=1cm]
    \def\w{6.6}
    % 填充（先画，避免遮住分隔线）
    \fill[blue!10] (0,3.0) rectangle (\w,4.0);
    \fill[blue!10] (0,2.0) rectangle (\w,3.0);
    \fill[black!4] (0,0.0) rectangle (\w,0.5);
    % 外框与分隔线
    \draw[draw=blue!45!black] (0,0) rectangle (\w,4.0);
    \draw[draw=blue!35!black] (0,0.5) -- (\w,0.5);
    \draw[draw=blue!35!black] (0,1.0) -- (\w,1.0);
    \draw[draw=blue!35!black] (0,1.5) -- (\w,1.5);
    \draw[draw=blue!35!black] (0,2.0) -- (\w,2.0);
    \draw[draw=blue!35!black] (0,3.0) -- (\w,3.0);
    % 字段文字
    \node at (\w/2,3.5) {\code{ra}：返回地址（\code{sd ra,24(sp)}）};
    \node at (\w/2,2.5) {\code{s0}：旧的 frame pointer（\code{sd s0,16(sp)}）};
    \node at (\w/2,1.75) {\code{n}：参数（\code{sw a0,-20(s0)}）};
    \node at (\w/2,1.25) {\code{i}：循环计数器（\code{-24(s0)}）};
    \node at (\w/2,0.75) {\code{sum}：累加结果（\code{-28(s0)}）};
    \node at (\w/2,0.25) {\footnotesize 对齐填充};
    % 指针标注
    \draw[-{Stealth[length=2mm]},thick,draw=blue!55!black] (-0.15,4.0) -- (0,4.0);
    \node[anchor=east,font=\footnotesize] at (-0.2,4.0) {\code{s0 = sp+32}（帧顶）};
    \draw[-{Stealth[length=2mm]},thick,draw=blue!55!black] (-0.15,0) -- (0,0);
    \node[anchor=east,font=\footnotesize] at (-0.2,0) {\code{sp}（帧底）};
  \end{tikzpicture}
  \caption{\code{calc} 的 32 字节栈帧（\code{addi sp,sp,-32}）}
\end{figure}

\subsection{手写汇编}

手写 \code{riscv/manual.s} 独立遵守同一约定：序言建立栈帧并保存 \code{ra}/\code{s0}，循环条件用 \code{blt} 在 \code{i > n} 时退出，奇偶判断用 \code{remw} 加 \code{beqz}，偶数项用 \code{slliw} 左移一位代替乘 2，累加用 \code{addw}，循环尾递增 \code{i} 后无条件跳回条件块。控制流标签与第 4 节的 LLVM 基本块一一对应。

\begin{lstlisting}
calc:
    addi sp, sp, -32
    sd   ra, 24(sp)
    sd   s0, 16(sp)
    addi s0, sp, 32
    sw   a0, -20(s0)        # n = 入参
    li   t0, 1
    sw   t0, -24(s0)        # i = 1
    sw   zero, -28(s0)      # sum = 0
.Lwhile_cond:
    lw   t0, -24(s0)        # i
    lw   t1, -20(s0)        # n
    blt  t1, t0, .Lwhile_end    # if (n < i) 退出循环
.Lif_cond:
    li   t1, 2
    remw t2, t0, t1         # i % 2
    beqz t2, .Lif_then
    j    .Lif_else
.Lif_then:
    lw   t0, -28(s0)        # sum
    lw   t1, -24(s0)        # i
    slliw t1, t1, 1         # 2 * i
    addw t0, t0, t1
    sw   t0, -28(s0)
    j    .Lif_end
.Lif_else:
    lw   t0, -28(s0)
    lw   t1, -24(s0)
    addw t0, t0, t1
    sw   t0, -28(s0)
.Lif_end:
    lw   t0, -24(s0)
    addiw t0, t0, 1         # i + 1
    sw   t0, -24(s0)
    j    .Lwhile_cond       # while 的回边
.Lwhile_end:
    lw   a0, -28(s0)        # 返回值放入 a0
    ld   s0, 16(sp)
    ld   ra, 24(sp)
    addi sp, sp, 32
    ret
\end{lstlisting}

\code{main} 只保留 16 字节栈帧用于存放 \code{ra}，连续调用 \code{getint}、\code{calc}、\code{putint}，三者之间的整数传递全部通过 \code{a0} 完成，即参数与返回值共用一个寄存器。

\subsection{链接与运行}

手写汇编由 \code{riscv64-unknown-elf-gcc} 汇编，并静态链接课程已有的 SysY runtime 归档 \code{libsysy_riscv.a}，沿用课程 \code{toolchain.conf} 与 \code{asm2bin.sh} 的配置：\code{-static -mcmodel=medany -Wl,--no-relax,-Ttext=0x90000000}。产物 \code{generated/manual_riscv} 是 RISC-V ELF 可执行文件，入口地址 \code{0x90000044}，由 \code{qemu-riscv64} 用户态模拟执行。runtime 的析构函数会向 stderr 打印 \code{TOTAL: 0H-0M-0S-0us} 计时摘要，测试脚本只比较标准输出，因此该摘要不影响数值判定。

\section{运行验证}

\begin{table}[H]
  \centering\small
  \caption{六组输入在三种实现上的输出比较}
  \label{tab:validation}
  \begin{tabular}{P{1.45cm}P{2.4cm}P{2.4cm}P{4.2cm}>{\centering\arraybackslash}p{1.6cm}}
    \toprule
    输入 $n$ & Reference & 手写 LLVM & 手写 RISC-V/QEMU & 结果 \\
    \midrule
    1 & 1 & 1 & 1 & PASS \\
    2 & 5 & 5 & 5 & PASS \\
    3 & 8 & 8 & 8 & PASS \\
    4 & 16 & 16 & 16 & PASS \\
    5 & 21 & 21 & 21 & PASS \\
    10 & 85 & 85 & 85 & PASS \\
    \bottomrule
  \end{tabular}
\end{table}

表~\ref{tab:validation} 的数字来自脚本 \code{scripts/09\_final\_validation.sh} 的实际输出。测试框架是 \code{source/test\_reference.sh}：它把六组 \code{输入 期望值} 硬编码为一个列表，逐组把输入喂给被测可执行文件，再与\emph{期望值}比较，而不是与另外两个实现的输出比较。这样三路一致但集体出错的情况同样会被检出，只有三者都等于 Stage 1 固定的 Ground Truth 才算通过。

各阶段结果分别留证于 \code{outputs/} 下的 \code{reference-test.txt}、\code{manual-llvm-test.txt} 与 \code{manual-riscv-test.txt}，汇总表为 \code{final-validation.md}，完整命令输出见 \code{logs/09_final_validation.log}。其中 RISC-V 一列由 \code{qemu-riscv64} 执行，课程 runtime 会在每次运行结束后向标准错误输出 \code{TOTAL: 0H-0M-0S-0us} 计时摘要；比较脚本只取标准输出，因此该摘要保留在日志中而不参与判定。

三份实现的栈帧、寄存器使用和指令选择各不相同，结果却完全一致，可佐证两套手写实现与源程序语义等价。

\section{优化观察：-O0 与 -O2}

对照产物由同一脚本生成：\code{exploration/O0.ll}、\code{O2.ll}（LLVM IR），\code{O0.s}、\code{O2.s}（x86-64），\code{O0\_riscv.s}、\code{O2\_riscv.s}（RV64），三份 unified diff 保存在 \code{outputs/optimization-*.diff}。以下只描述当前程序真实发生的变化。

\begin{table}[H]
  \centering\small
  \caption{LLVM IR 指令数量与结构变化}
  \begin{tabularx}{\textwidth}{P{3.2cm}P{3.6cm}Y}
    \toprule
    指标 & -O0 & -O2 \\
    \midrule
    \code{alloca} & 6 & 0 \\
    \code{load} & 11 & 0 \\
    \code{store} & 9 & 0 \\
    \code{phi} & 0 & 32 \\
    局部变量表示 & 栈槽 + 内存读写 & SSA 临时值 + 循环 \code{phi} \\
    奇偶判断 & \code{srem} + \code{icmp eq} + \code{br} & \code{and}/\code{xor}/\code{shl} 位运算 \\
    \code{main} 对 \code{calc} 的调用 & 存在 \code{call @calc} & 已内联；\code{@calc} 定义因外部可见而保留 \\
    \bottomrule
  \end{tabularx}
\end{table}

\subsection{从内存槽到 SSA}

\code{-O0} 把每个局部变量放进取址栈槽，任何读写都要经过 \code{load}/\code{store}。\code{-O2} 在 IR 上执行 mem2reg 一类变换后，可变变量成为 SSA 值；循环中跨迭代传递的量则用 \code{phi} 节点表达。本例中 \code{alloca/load/store} 全部归零、\code{phi} 从 0 增至 32。这一步先改变的是\emph{值的表示方式}而非指令条数；寄存器化之后，算术化简和向量化才有作用对象。

\subsection{奇偶判断被改写成位运算}

源程序用 \texttt{i \% 2 == 0} 选择分支。\code{-O2} 识别出这是奇偶测试，不再依赖取模。x86-64 的 IR 中可以看到 \code{and}/\code{xor}/\code{shl} 的组合，RV64 的 \code{calc} 则压缩为：

\begin{lstlisting}
calc:
    li   a1, 0
    blez a0, .LBB0_3
    li   a2, 1
.LBB0_2:
    mv   a3, a2
    not  a2, a2
    andi a2, a2, 1
    sllw a2, a3, a2        # i << (i 为偶数 ? 1 : 0)
    addw a1, a1, a2        # sum += 该项
    addiw a2, a3, 1
    blt  a3, a0, .LBB0_2
.LBB0_3:
    mv   a0, a1
    ret
\end{lstlisting}

\code{not}/\code{andi} 取出 $i$ 的奇偶位再取反，得到“$i$ 为偶数时为 1，否则为 0”的移位量；\code{sllw a2, a3, a2} 做变长左移，同时覆盖了原程序的两个分支：偶数时左移 1 位即乘 2，奇数时左移 0 位即原值。整个 \code{if/else} 因此消失，只剩一个循环。此外 \code{calc} 在 \code{-O2} 下没有建立任何栈帧（无 \code{sp} 调整、无 \code{ra} 保存），因为优化后它已不含函数调用，成为叶子函数。

\subsection{两种目标的低层差异}

尽管 IR 是目标无关的，同一份 \code{-O2} IR 在两种后端上继续下降的结果并不相同。

x86-64 带 SSE2，且本程序的循环次数可判定，因此 loop vectorizer 把它改写为四 lane 运算：\code{O2.ll} 中出现 \code{<4 x i32>} 类型的 \code{phi} 与 \code{shl}，以及 \code{call i32 @llvm.vector.reduce.add.v4i32}，汇编中可见 XMM 寄存器和 \code{paddd}。由于 \code{calc} 不是 4 的整数倍长度也能正确退出，向量化版本还额外生成了边界与尾部路径。

RV64 目标为 \code{rv64gc}，不含 RVV 向量扩展，因此 \code{O2\_riscv.s} 仍是标量循环，没有向量寄存器指令；优化所做的只是内联与算术化简。两种目标的共同点是 \code{main} 都不再调用 \code{calc}，即函数内联在两条路径上都发生了。

\begin{table}[H]
  \centering\small
  \caption{\code{-O0} 与 \code{-O2} 汇编规模变化}
  \begin{tabular}{P{3.4cm}rrP{5.4cm}}
    \toprule
    汇编文本行数 & -O0 & -O2 & 变化来源 \\
    \midrule
    x86-64（\code{.s}） & 87 & 385 & 向量化循环体展开与 \code{calc} 内联 \\
    RV64（\code{.s}） & 91 & 56 & 消除栈帧读写、分支退化为变长移位 \\
    \bottomrule
  \end{tabular}
\end{table}

两行行数朝相反方向变化，说明\emph{汇编行数不能作为优化效果的度量}：x86-64 因为向量化展开了循环而变长，RV64 因为化简而变短。本实验只观察 IR 与汇编，没有测量运行耗时、代码体积或能耗，因此不据此声称任何量化的性能提升。

\section{MLIR / AscendNPU VecAdd}

MLIR 以 operation/value/block/region 和 dialect 描述不同抽象层。多层表示的作用是让每一阶段保留它需要的信息：张量形状和算子边界便于做 fusion、tiling 与调度；进一步 lowering 后，内存位置、数据搬运和目标硬件操作才逐步明确。若过早把算子摊平成低层指令，后续优化就难直接利用张量和算子语义。LLVM IR 可作为后续通用低层表示，但它不替代前面针对张量和 NPU 硬件的抽象。\cite{mlir-lang,mlir-rationale,mlir-llvm,ascend-framework}

\begin{figure}[H]
  \centering
  \begin{tikzpicture}[
    stage/.style={draw=blue!45!black,fill=blue!4,rounded corners=2pt,
      align=center,text width=2.65cm,minimum height=0.86cm,font=\small},
    flow/.style={-{Stealth[length=2mm]},thick,draw=blue!55!black},
    node distance=0.5cm
  ]
    \node[stage] (tensor) {Arith / Tensor\\元素类型与 shape};
    \node[stage,right=of tensor] (hfusion) {Linalg / HFusion\\算子、融合与 tiling};
    \node[stage,right=of hfusion] (hivm) {HIVM\\GM / UB、load / vadd / store};
    \node[stage,right=of hivm] (lower) {HACC / LLVM\\后续入口与目标 lowering};
    \draw[flow] (tensor) -- (hfusion);
    \draw[flow] (hfusion) -- (hivm);
    \draw[flow] (hivm) -- (lower);
  \end{tikzpicture}
  \caption{VecAdd 的概念 lowering 主线；本次逐行核对与实跑的实物位于 HIVM 层}
  \label{fig:lowering}
\end{figure}

\begin{table}[H]
  \centering\small
  \caption{VecAdd lowering 各层表达的信息}
  \begin{tabularx}{\textwidth}{P{2.5cm}P{5.2cm}Y}
    \toprule
    层 & VecAdd 表达 & 关注点 \\
    \midrule
    Arith/Tensor & 张量形状、元素类型和逐元素加 & 算法语义，未指定 NPU 地址空间 \\
    Linalg/HFusion & 张量算子及融合、tiling & 算子组合和执行计划 \\
    HIVM & GM/UB、load、vadd、store & 显式表示数据搬运和 NPU 运算位置 \\
    HACC & entry 与 HOST/DEVICE 属性 & 主机/设备边界与 kernel 入口 \\
    LLVM IR/后端 & 更低层通用表示和目标代码 & 后续 lowering；此处不是本次实测 dump \\
    \bottomrule
  \end{tabularx}
\end{table}

官方 HIVM VecAdd 文件 \code{mlir/sources/vecadd_hivm.mlir} 中，16 个 \code{i16} 元素的 A、B、C 参数位于 GM；两次 \code{hivm.hir.load} 把输入搬到 UB，\code{hivm.hir.vadd} 在 UB 上计算，\code{hivm.hir.store} 再写回 GM。这里的数据不是只标注为“向量值”，还带有存储位置和搬运路径。

\begin{lstlisting}
func.func @add(%arg0: memref<16xi16, #hivm.address_space<gm>>,
               %arg1: memref<16xi16, #hivm.address_space<gm>>,
               %arg2: memref<16xi16, #hivm.address_space<gm>>)
    attributes {hacc.entry, hacc.function_kind = #hacc.function_kind<DEVICE>} {
  %alloc   = memref.alloc() : memref<16xi16, #hivm.address_space<ub>>
  hivm.hir.load ins(%arg0 : ...<gm>) outs(%alloc   : ...<ub>)   // GM -> UB
  %alloc_0 = memref.alloc() : memref<16xi16, #hivm.address_space<ub>>
  hivm.hir.load ins(%arg1 : ...<gm>) outs(%alloc_0 : ...<ub>)   // GM -> UB
  %alloc_1 = memref.alloc() : memref<16xi16, #hivm.address_space<ub>>
  hivm.hir.vadd ins(%alloc, %alloc_0 : ...<ub>, ...<ub>) outs(%alloc_1 : ...<ub>)
  hivm.hir.store ins(%alloc_1 : ...<ub>) outs(%arg2 : ...<gm>)  // UB -> GM
  return
}
\end{lstlisting}

与高层表示相比，这里的 \code{memref.alloc} 都显式标注 \code{<ub>} 地址空间，\code{hivm.hir.load}/\code{store} 把“从哪搬到哪”写成显式指令。NPU 编程与普通标量程序的区别在于：计算必须在片上存储（UB）上进行，数据要先从全局内存（GM）搬入、算完再搬回。这些原本隐含的硬件执行策略由此被写入 IR。\cite{ascend-framework,ascend-vecadd}

\subsection{PlanMemory 上游用例分析}

PlanMemory 根据 buffer 生命周期和硬件约束规划片上地址。它把局部 \code{memref.alloc} 改写为带偏移的 \code{hivm.hir.pointer_cast}。

Ascend 上游测试文件 \code{plan-memory.mlir} 中的 \texttt{test\allowbreak\_mem\allowbreak\_allocate\allowbreak\_basic} 用例包含四个大小为 $16\times16\times16\times2=8192$ 字节的 UB 临时 buffer。其上游 RUN 命令与关键断言为：

\begin{lstlisting}
// RUN: bishengir-opt %s -hacc-append-device-spec=target=Ascend910B1 \
//        -hivm-plan-memory -split-input-file -verify-diagnostics | FileCheck %s
// CHECK-LABEL: func.func @test_mem_allocate_basic
// CHECK-NOT: memref.alloc()
// CHECK: %[[CONST0:.*]] = arith.constant 0 : i64
// CHECK: {{.*}} = hivm.hir.pointer_cast(%[[CONST0]])
\end{lstlisting}
\code{RUN} 行中的 \code{-hacc-append-device-spec=target=Ascend910B1} 是必需项：PlanMemory 要先知道目标芯片规格，才能应用相应的对齐与容量约束。该文件还有第二条 \code{RUN}，以 \code{=largest-first} 选择另一种分配策略并用 \code{CHECK-REORDER} 前缀检查。

\subsection{实测：PlanMemory 把 alloc 改写为 pointer\_cast}

本实验从源码构建了上游 \code{bishengir-opt}（见本节的"构建与验证方式"），并逐字执行上面两条 \code{RUN} 命令，把输出交给\emph{上游自己的} \code{CHECK} 断言做 FileCheck：默认策略与 \code{largest-first} 两条均为 PASS。

对用例 \code{test\_mem\_allocate\_basic} 单独保存的 Before/After 如下。输入含四个 16$\times$16$\times$16$\times$\code{f16} 的 UB 临时 buffer，共 32 KiB。

\begin{table}[H]
  \centering\small
  \caption{PlanMemory 在该用例上的实测变化}
  \begin{tabular}{P{4.4cm}rr}
    \toprule
    指令计数 & Before & After \\
    \midrule
    \code{memref.alloc()} & 4 & 0 \\
    \code{hivm.hir.pointer\_cast} & 0 & 4 \\
    \code{hivm.hir.load}/\code{vadd}/\code{store} & 5 & 5 \\
    \bottomrule
  \end{tabular}
\end{table}

四个 \code{memref.alloc()} 被逐个改写为引用同一常量 \code{\%c0\_i64 = arith.constant 0 : i64} 的 \code{pointer\_cast}，偏移均为 0；而计算指令一条未改。这与此前只能引用文档的推断一致，但现在有本机产出的 Before/After 为证（\code{mlir/outputs/plan\_memory\_*.mlir}）。

Pass 同时把目标芯片规格写入了模块属性，这是"硬件细节何时进入 IR"的具体形态：

\begin{lstlisting}
module attributes {
  dlti.target_system_spec = #dlti.target_system_spec<"NPU" :
    #hacc.target_device_spec<
      #dlti.dl_entry<"UB_SIZE", 1572864 : i32>,      // 1.5 MiB
      #dlti.dl_entry<"L1_SIZE", 4194304 : i32>,      // 4 MiB
      #dlti.dl_entry<"L0C_SIZE", 1048576 : i32>,
      #dlti.dl_entry<"UB_ALIGN_SIZE", 256 : i32>,
      #dlti.dl_entry<"AI_CORE_COUNT", 24 : i32>,
      #dlti.dl_entry<"VECTOR_CORE_COUNT", 48 : i32>, ...>>,
  hacc.target = #hacc.target<"Ascend910B1">
}
\end{lstlisting}

关于偏移为 0：该用例四个 buffer 共 32 KiB，远小于芯片 UB 容量（1572864 字节），因此都落在偏移 0 处。不能把单个用例的 offset-0 期望推广到任意程序——一般情况仍需满足生命周期、别名、对齐和硬件复用规则。\cite{ascend-planmemory,ascend-passes,ascend-planmemory-test}

\subsection{实测：InjectSync 注入流水线同步}

上游 \code{inject-sync.mlir} 中的 \code{test\_mem\_injct\_sync\_basic} 用例，源程序只有 \code{load}、\code{vadd}、\code{store} 三条计算指令，不含任何同步。逐字执行其 \code{RUN} 命令（\code{-hivm-inject-sync -hivm-lower-multi-buffer-counter}）后，注入结果如下：

\begin{lstlisting}
hivm.hir.load ins(%arg0 : ...<gm>) outs(%0 : ...<ub>)
hivm.hir.set_flag[<PIPE_MTE2>, <PIPE_V>, <EVENT_ID0>]
hivm.hir.wait_flag[<PIPE_MTE2>, <PIPE_V>, <EVENT_ID0>]
hivm.hir.vadd ins(%0, %2 : ...) outs(%0 : ...)
hivm.hir.set_flag[<PIPE_V>, <PIPE_MTE3>, <EVENT_ID0>]
hivm.hir.wait_flag[<PIPE_V>, <PIPE_MTE3>, <EVENT_ID0>]
hivm.hir.store ins(%0 : ...<ub>) outs(%arg1 : ...<gm>)
hivm.hir.pipe_barrier[<PIPE_ALL>]
\end{lstlisting}

这段输出把"NPU 为什么需要显式同步"讲得很具体：数据搬运走 MTE2/MTE3 流水，计算走 V（向量）流水，二者是\emph{不同的硬件流水线}，源代码里的先后书写顺序并不保证执行完成。Pass 因此插入 MTE2$\to$V（搬入后才能算）与 V$\to$MTE3（算完才能写回）两组 flag，并在函数末尾补一条 \code{PIPE\_ALL} barrier。计数上，Before 为 0 条同步指令，After 为 2 条 \code{set\_flag}、2 条 \code{wait\_flag} 和 1 条 \code{pipe\_barrier}（\code{mlir/outputs/sync\_*.mlir}）。

\subsection{构建与验证方式}

本节的 Ascend 结果全部来自本机运行，验证方式是执行上游测试文件自带的 \code{RUN} 命令，再用上游的 \code{CHECK} 断言做 FileCheck——即用上游自己的验收标准检验本机构建。

工具链不在系统内：系统 CMake 缺失、Ninja 为 1.10.1（官方要求至少 1.12），因此 CMake 4.4.4 与 Ninja 1.13.2 由 pip 安装在仓库外的用户级 venv 中，AscendNPU-IR 及其 LLVM 19.1.7 子模块按固定 commit 检出于 \code{third\_party/}，未改动系统环境。构建产物自报版本为：

\begin{lstlisting}
bishengir-opt 2.0.0 (gitcode.com/Ascend/ascendnpu-ir 2ac1b376da1d)
llvm 19.1.7 06913395ef8a   Release build
\end{lstlisting}

构建过程中对上游源码打了一处必要的链接补丁：\code{BiShengIRHIVMCLib} 与 \code{BiShengIRHIVMCLibA3/A5} 之间存在真实的静态库循环（A3/A5 调用前者的 \code{setSanitizerAddrArgName}，前者又调用后者的 \code{runHIVMCCompileA3/A5}），静态库单遍扫描无法解析。补丁用 CMake 的 \code{\$<LINK\_GROUP:RESCAN,...>} 让链接器对这组库重复扫描。它只影响 \code{bishengir-opt} 的链接参数，不改变任何 Pass 的行为——三个上游 \code{RUN} 命令连同其断言全部通过，即为此提供了佐证。补丁全文见 \code{scripts/patches/}。

\subsection{边界}

第 8 节的实测范围止于 IR 层：本机运行了 Pass，但没有安装完整 CANN、没有配置 Ascend 驱动、没有生成或运行 kernel \code{.o}，也没有在真实 NPU 上执行。已确认这些对"观察 MLIR 中端 lowering"不是必需项——CANN 只在编译 kernel 并上板执行时才需要。

另需区分：图~\ref{fig:lowering} 中的 Arith/Tensor 与 HFusion 两层是依据官方架构资料整理的概念主线，本次逐行核对与实跑的实物位于 HIVM 层。\cite{ascend-build,ascend-readme}

\clearpage
\section{实验可复现性}

全部实验由 \code{scripts/} 下的脚本驱动，每个阶段一条命令即可重放，输出同时写入 \code{logs/} 与 \code{records/}。自动化与手写产物分属不同目录：\code{generated/} 只放工具生成结果，\code{llvm/} 与 \code{riscv/} 放人工编写结果，避免把编译器输出当成交付内容。

\begin{table}[H]
  \centering\small
  \caption{各阶段复现入口}
  \begin{tabularx}{\textwidth}{P{5.4cm}Y}
    \toprule
    脚本 & 作用 \\
    \midrule
    \code{00_environment_check.sh} & 记录系统与工具版本，作为环境快照 \\
    \code{01_reference.sh} & 编译宿主机 Ground Truth 并跑六组测试 \\
    \code{02_pipeline.sh} & 生成 \code{.i / .ll / .s / .o} 并显式链接 \\
    \code{03_llvm_analysis.sh}、\code{03_llvm_validation.sh} & 抽取 IR 做 CFG 分析；验证手写 \code{manual.ll} \\
    \code{04_riscv_reference.sh}、\code{05_riscv_validation.sh} & 生成 RV64 参考汇编；汇编链接手写 \code{manual.s} 并在 QEMU 运行 \\
    \code{06_optimization.sh} & 生成 O0/O2 三种目标的对照与 diff \\
    \code{08_mlir_exploration.sh} & 记录进阶部分的只读环境检查与范围判定 \\
    \code{09_final_validation.sh} & 三路实现统一比对并输出总表 \\
    \code{10_build_report.sh} & 用 Tectonic 编译本报告 \\
    \midrule
    \code{11_ascend_configure.sh} & 配置 AscendNPU-IR（只配置不编译），校验 \code{bishengir-opt} / \code{hivmc} target \\
    \code{12_ascend_build.sh} & 编译 \code{bishengir-opt} 单个 target \\
    \code{13_ascend_pass_runs.sh} & 逐字执行上游 RUN 命令并用其 CHECK 做 FileCheck；保存两组 Before/After \\
    \bottomrule
  \end{tabularx}
\end{table}

第 11–13 个脚本针对进阶部分，依赖仓库外的源码与工具链（\code{third\_party/} 与 \code{.venv-ascendnpu}，二者不进版本管理；获取方式见第 11 节的记录）。前 10 个脚本只依赖课程环境。

统一执行形式为 \code{bash scripts/<name>.sh 2>&1 | tee logs/<name>.log}。脚本以 \code{set -euo pipefail} 开头，任一步失败即中止，因此 \code{logs/} 中的记录不会掩盖中途错误。实验前需先 \code{source ../env.sh} 以取得 QEMU 与 RISC-V 工具链路径；\code{09_final_validation.sh} 已内置该步骤。

\section{结论}

本文用一个包含函数、循环、分支、取模和运行库调用的 SysY 程序，完成了从源程序到可执行文件的全部环节，并在此基础上独立完成了两套手写实现。按引言提出的四个问题，可以得到以下结论。

\textbf{四个阶段的分工。} 预处理只做文本替换，产物仍是源文本；前端把它变成 LLVM IR，其间经过词法、语法和语义分析；汇编器把汇编变成可重定位目标文件，此时地址未定、外部符号未解；链接器才把定义与引用配对并定下最终地址。本实验的 \code{example.o} 是 \code{ET\_REL}、入口地址为 0、含两个未定义符号，而链接后的文件六组测试全部通过，这一对比即为四者职责的直接证据。

\textbf{IR 到源代码的映射。} \code{-O0} 下 IR 与源程序接近一一对应：\code{alloca/load/store} 表示变量，\code{icmp}/\code{br} 表示比较与转移，\code{define/call/ret} 表示函数关系。\code{while} 和 \code{if} 都没有专用指令，它们由基本块、条件计算和 \code{br} 组合而成；循环由 \code{if.end} 跳回 \code{while.cond} 的回边体现。正因为这一映射是结构性的，才能在阅读自动生成结果之后独立重写出一份语义等价、命名更清晰的 \code{manual.ll}。

\textbf{RISC-V 的落地方式。} 参数与返回值共用 \code{a0}，\code{calc} 用 32 字节栈帧保存 \code{ra}、\code{s0} 和三个局部量，控制流落为比较加条件跳转的标签结构。手写版与 Clang 自动版采用了\emph{相同的栈帧布局}，但在取模和 \code{main} 的帧大小上取舍不同，可见 ABI 约束是刚性的，具体实现仍有空间。

\textbf{优化与抽象层次。} \code{-O2} 先改变值的表示（栈槽变 SSA，出现 32 个 \code{phi}），再重写算术（奇偶判断退化为变长左移，整个 \code{if/else} 消失），最后才由后端按目标能力决定是否向量化：x86-64 生成四 lane 运算和 \code{paddd}，RV64 因缺少 RVV 仍是标量循环。两种目标的汇编行数一增一减（87$\to$385 与 91$\to$56），说明行数不能作为优化效果的度量。MLIR 部分从另一个方向回答了分层问题，且本轮已有实测支撑：PlanMemory 把四个 \code{memref.alloc()} 改写为带偏移的 \code{pointer\_cast} 并注入含 UB 容量与对齐的芯片规格，而计算指令一条未改；InjectSync 在只有 load/vadd/store 的用例上注入 MTE2$\to$V 与 V$\to$MTE3 两组同步和一条 \code{PIPE\_ALL} barrier。两者共同说明，硬件搬运与流水线语义是被逐步、显式地写进 IR 的。

\textbf{本实验的边界。} 八节结论均基于本机真实产物，可用第 9 节的脚本逐阶段复现。第 8 节的 Ascend 部分虽从源码构建了 \code{bishengir-opt} 并实跑了两组 Pass，但范围止于 IR 层：没有安装 CANN、没有配置驱动、没有生成或运行 kernel \code{.o}、没有在真实 NPU 上执行。这些经确认不是观察 MLIR 中端 lowering 的必需项。此外，lowering 主线中的 Arith/Tensor 与 HFusion 两层仍属依据官方资料整理的概念说明，实物只到 HIVM 层。

\begin{thebibliography}{99}
\footnotesize
\setlength{\itemsep}{0pt}
\setlength{\parsep}{0pt}
\setlength{\topsep}{0.2em}
\bibitem{mlir-lang} LLVM Project. \emph{MLIR Language Reference}. \url{https://mlir.llvm.org/docs/LangRef/}（访问日期：2026-10-07）.
\bibitem{mlir-rationale} LLVM Project. \emph{MLIR Rationale}. \url{https://mlir.llvm.org/docs/Rationale/Rationale/}（访问日期：2026-10-07）.
\bibitem{mlir-llvm} LLVM Project. \emph{Targeting LLVM IR from MLIR}. \url{https://mlir.llvm.org/docs/TargetLLVMIR/}（访问日期：2026-10-07）.
\bibitem{ascend-framework} Ascend. \emph{AscendNPU-IR Framework Integration}. \url{https://github.com/Ascend/AscendNPU-IR/blob/master/docs/source/en/developer_guide/conversion/framework_interface.md}（访问日期：2026-10-07）.
\bibitem{ascend-interface} Ascend. \emph{AscendNPU-IR Interface API}. \url{https://github.com/Ascend/AscendNPU-IR/blob/master/docs/source/en/developer_guide/conversion/interface_api.md}（访问日期：2026-10-07）.
\bibitem{ascend-vecadd} Ascend. \emph{HIVM VecAdd add.mlir}. \url{https://github.com/Ascend/AscendNPU-IR/blob/master/bishengir/test/Integration/HIVM/VecAdd/add.mlir}（访问日期：2026-10-07）.
\bibitem{ascend-readme} Ascend. \emph{HIVM VecAdd README and host wrapper}. \url{https://github.com/Ascend/AscendNPU-IR/blob/master/bishengir/test/Integration/HIVM/VecAdd/README.md}（访问日期：2026-10-07）.
\bibitem{ascend-planmemory} Ascend. \emph{Plan Memory}. \url{https://ascendnpu-ir.gitcode.com/en/developer_guide/features/PlanMemory/PlanMemory.html}（访问日期：2026-10-07）.
\bibitem{ascend-passes} Ascend. \emph{HIVM Passes}. \url{https://ascendnpu-ir.gitcode.com/en/developer_guide/passes/hivm_passes.html}（访问日期：2026-10-07）.
\bibitem{ascend-planmemory-test} Ascend. \emph{PlanMemory Test Case}. \url{https://github.com/Ascend/AscendNPU-IR/blob/master/bishengir/test/Dialect/HIVM/plan-memory.mlir}（访问日期：2026-10-07）.
\bibitem{ascend-build} Ascend. \emph{AscendNPU-IR Build and Installation}. \url{https://github.com/Ascend/AscendNPU-IR/blob/master/docs/source/en/introduction/quick_start/installing_guide.md}（访问日期：2026-10-07）.
\bibitem{course-docs} Compiler-Labs. 编译原理作业要求与执行流程（本地课程文档，访问日期：2026-10-07）.
\end{thebibliography}

\end{document}
